The defining conditions say exactly that $Q$ is the image of $P$ under rotation about $A$ through $\theta$. A rotation maps a line onto a line, so the locus is $\boxed{\text{the rotated image of }\ell}$, including every point of that image because the rotation is invertible.
For an unoriented angle with $0<\theta<180^\circ$, both rotations by $\theta$ and $-\theta$ are allowed. The locus is the union of their two image lines, which may coincide in special cases. At $0^\circ$ or $180^\circ$ the two rotations agree, giving one line. If $A$ lies on $\ell$, its image $A$ is included as the limiting and degenerate case.
\begin{asy}
texpreamble("\usepackage{amsmath,amssymb}");
size(240);
defaultpen(fontsize(9pt));
pair A=(0,0),P=(1,1),Q=rotate(40)*P;draw((-2,1)--(2,1));draw(rotate(40)*(-2,1)--rotate(40)*(2,1),blue);draw(rotate(-40)*(-2,1)--rotate(-40)*(2,1),gray+dashed);draw(P--A--Q);dot(A);dot(P);dot(Q);label("$A$",A,(0,-1));label("$P$",P,(1,0));label("$Q$",Q,(0,1));label("$\ell$",(2,1),(1,0));label("$+\theta$",(0.6,.25));
\end{asy}
